\chapter*{\centering Abstract}
\begin{singlespace}
\quad The primary objective of this work is to identify, address, and validate critical
security and privacy gaps in DAuth, a lightweight decoupled distributed authentication
scheme for multihop IoT networks.
Two significant vulnerabilities were identified.
First, the real device identity is transmitted in plaintext over the public
channel during every authentication session.
This exposure enables any passive adversary to track devices across sessions,
correlate repeated authentication attempts, and infer network behaviour, which
directly violates modern privacy requirements.
Second, when a key-exchange packet is lost in transit, the device and authentication
server hold divergent session states, causing all subsequent authentication attempts
to fail and forcing the device through a costly full re-enrolment.

To address both vulnerabilities, this work proposes an extension of DAuth that
preserves device anonymity and tolerates desynchronisation.
Per-session pseudonym rotation replaces the real device identity in all
open-channel communications with a rotating pseudonym that is refreshed after each
successful key exchange, so that successive pseudonyms are computationally unlinkable.
A dual-state mechanism maintained at the authentication server retains both the
current and the immediately preceding session state, enabling seamless recovery
from a single packet-loss event without re-enrolment.
Formal verification under a standard adversary model confirms authentication
correspondence, session key secrecy, forward secrecy, anonymity preservation, and
offline guessing resistance.
COOJA\slash Contiki-NG simulations on a 100-node RPL network show a 42.8\,\% reduction
in per-device energy over LAAKA and 32.4\,\% over Zhou et al., with consistent gains
across all network sizes and authentication server pool configurations.
Hardware evaluation on Raspberry Pi~4B confirms 52.5\,\% and 14.2\,\% energy
reductions over LAAKA and Zhou et al., respectively; DAuth achieves the lowest
hardware overhead as it carries no pseudonym operations, while the proposed scheme
remains well below both LAAKA and Zhou et al.
\end{singlespace}
